\documentclass[10pt,a4paper]{amsart}
\usepackage[margin=19mm]{geometry}
\usepackage{amsmath,amssymb,mathtools}
\usepackage[scr=rsfs,cal=euler]{mathalfa}
\usepackage{xfrac}
\usepackage[unicode,hidelinks,bookmarks=false]{hyperref}
\newcommand{\ie}{i.e.\ }
\newcommand{\cf}{cf.\ }
\newcommand{\resp}{respectively\ }
\newcommand{\Subsection}{\subsection}
\providecommand{\nobreakdash}{\nobreak\hbox{-}}
\setlength{\emergencystretch}{2em}
\multlinegap0pt
\usepackage{bm}
\usepackage{bbm}
\usepackage{dsfont}
\usepackage{xspace}
\usepackage{cancel}
\usepackage{mathtools}
\usepackage{amsthm}
\makeatletter
\@ifundefined{lemma}{\newtheorem{lemma}{Lemma}}{}
\@ifundefined{prop}{\newtheorem{prop}{Proposition}}{}
\makeatother
\title[Stability of intrinsic localized modes on the lattice with c]{Stability of intrinsic localized modes on the lattice with competing power nonlinearities}
\author{Georgy L. Alfimov, Pavel A. Korchagin, Dmitry E. Pelinovsky}
\date{Source: arXiv:2511.12649v1 (CC BY 4.0)}
\begin{document}
\maketitle

\noindent\emph{Source and licence.} Stability of intrinsic localized modes on the lattice with competing power nonlinearities. Version arXiv:2511.12649v1 (CC BY 4.0), distributed under the Creative Commons Attribution 4.0 International licence, as recorded in the arXiv licence metadata for this exact version and shown on the abstract page https://arxiv.org/abs/2511.12649v1. Full rights notice: licenses/arxiv-2511.12649-CC-BY-4.0.txt.\par\smallskip

\noindent\textit{Editorial scope.} Counted unit: the Prop whose source label is \texttt{prop-existence} in the article (source0.tex lines 346--358, source label \texttt{prop-existence}), delivered together with the article's own complete original proof of it (source0.tex lines 360--408, one \texttt{proof} environment of 49 lines). No mathematics is written, repaired or re-derived: statement and proof are the article's own text line for line. Required context retained uncounted: Eq:u\_gen\_pq (lines 128--131); lem-1 (lines 261--264). Every definition, display and label of those lines is retained unchanged. Omitted from this package and not redistributed: 1--127, 132--260, 265--345, 359--359, 409--2871. Nothing inside a retained excerpt is omitted or altered: every retained body is delivered exactly as the article's lines are written, so no omission receipt of any kind accompanies this package. Citations: the retained excerpts carry no citation command at all, so no bibliography entry of the article is transcribed and no reference is redistributed. Reference adaptation: none was needed and none was made; every reference inside the delivered unit resolves inside it, so no unresolved reference or citation is produced. Printed slips kept literal: no printed word, symbol, hypothesis or number of the counted statement or of its proof was altered. Layout and class: the article's own class is \texttt{amsart} with packages including graphicx, amssymb, amsmath, amsthm, color, pstricks, hyperref, empheq, pgfplots, comment, booktabs, multirow, array, cancel; the wrapper is the lane's pinned \texttt{amsart} editorial wrapper at 10pt on A4, which restates the theorem-like environments the retained text needs and adds the article's own macro definitions that the retained text needs (none), with extra wrapper packages (bm, bbm, dsfont, xspace, cancel, mathtools). The delivered numbering is the unit's own. Assets: no retained span contains a figure environment, a graphics-inclusion command, a PDF-page inclusion, a table or a TikZ picture; no image file is copied into this package, the package declares no asset dependency and no image file is redistributed with it. Licence: version arXiv:2511.12649v1 of this article is distributed under the Creative Commons Attribution 4.0 International licence, as recorded in the arXiv licence metadata for this exact version and shown on the abstract page https://arxiv.org/abs/2511.12649v1. A full rights notice is delivered at licenses/arxiv-2511.12649-CC-BY-4.0.txt. Characters: every retained excerpt is pure ASCII, so the delivered engine, XeLaTeX with its default Latin Modern text font, reports no missing character.
\begin{gather}
    \varepsilon\left(u_{n+1}-2u_n+u_{n-1}\right)- u_n+\left|u_n\right|^{p-1} u_n-\gamma\left|u_n\right|^{q-1} u_n=0. 
    \label{Eq:u_gen_pq}
\end{gather}

\begin{lemma}
	\label{lem-1}
	For any $p,q$ with $q>p$, and $0<\gamma<\gamma_{{p,q}}$, we have  $f'(a) < 0$ and $f'(A) > 0$.
\end{lemma}

\begin{prop}
	\label{prop-existence}
	Let $p,q \in \mathbb{N}$ be fixed such that $2 \leq p < q$. 
	For every $\gamma \in (0,\gamma_{{p,q}})$, there exists $\varepsilon_0 > 0$ and $C_0 > 0$ such that for every $\varepsilon \in (0,\varepsilon_0)$, there exists 
	a solution ${\bf u} \in \ell^2(\mathbb{Z})$ of the difference 
	equation (\ref{Eq:u_gen_pq}) such that 
\begin{equation}
\label{bound-u}
	\| {\bf u} - {\bf u}^{(0)} \|_{\ell^2(\mathbb{Z})} \leq C_0 \varepsilon,
\end{equation}
	where ${\bf u}^{(0)} = (\ldots,0,0,\tilde{\bf u},0,0,\ldots)$ with  $\tilde{\bf u} \in \mathbb{R}^N$ of any length $N \geq 1$. Moreover, 
	the number of sign alternations in ${\bf u}$ is equal to the number of sign alternations in $\tilde{\bf u}$.
\end{prop}

\begin{proof}
	By writing ${\bf u} = {\bf u}^{(0)} + {\bf v}$ with real-valued ${\bf v} \in \ell^2(\mathbb{Z})$, we rewrite the difference 
	equation (\ref{Eq:u_gen_pq}) in the equivalent form:
\begin{equation}
-\varepsilon \left(v_{n+1}-2v_n+v_{n-1}\right) + f'(u_n^{(0)}) v_n = \varepsilon \left( u_{n+1}^{(0)} - 2 u_n^{(0)} + u_{n-1}^{(0)} \right) - g(u_n^{(0)},v_n),
\label{Eq:v}
\end{equation}	
where $f(u) = u - |u|^{p-1} u + \gamma |u|^{q-1} u$ and $g(u,v) =  f(u + v) - f(u) - f'(u) v$. Since $p,q \in \mathbb{N}$ and $2 \leq p < q$, the mapping $v \to g(u,v)$ is $C^1$ for a given $u \in \mathbb{C}$ and 
there is $C > 0$ such that 
$$
|g(u_n^{(0)},v_n)| \leq C|v_n|^2, \quad \forall n \in \mathbb{Z},
$$
as long as $\|{\bf v} \|_{\ell^2} \leq 1$. We can rewrite (\ref{Eq:v}) in the abstract form:
\begin{equation}
\label{eq:v-abstract}
\mathcal{L} {\bf v} = {\bf H}({\bf v}), 
\end{equation}
with 
\begin{align*}
\mathcal{L} v_n &= f'(u_n^{(0)}) v_n - \varepsilon \Delta_2 v_n, \\
H_n({\bf v}) &= \varepsilon \Delta_2 u_n^{(0)} - g(u_n^{(0)},v_n).
\end{align*}
The mapping $\ell^2(\mathbb{Z}) \ni {\bf v} \to {\bf H}({\bf v}) \in \ell^2(\mathbb{Z})$ is $C^1$ since $\ell^2(\mathbb{Z})$ is a Banach algebra with respect to multiplication and
$$
\| \Delta_2 {\bf u}^{(0)} \|_{\ell^2} \leq 4 \|{\bf  u}^{(0)} \|_{\ell^2}. 
$$
On the other hand, we have $f'(0) = 1$, $f'(\pm a) < 0$, $f'(\pm A) > 0$ by Lemma \ref{lem-1}. Hence, $\mathcal{L}$ is an invertible operator in $\ell^2(\mathbb{Z})$ for sufficiently small $\varepsilon > 0$ with a bounded inverse. By the implicit function theorem in $\ell^2(\mathbb{Z})$, there exist $\varepsilon_0 > 0$ and $C_0 > 0$ such that there exists a unique solution 
to (\ref{eq:v-abstract})  satisfying $\| {\bf v} \|_{\ell^2} \leq C_0 \varepsilon$ for every $\varepsilon \in (0,\varepsilon_0)$. This yields the first assertion and results in the bound (\ref{bound-u}). 

It remains to prove preservation of the sign alternation in ${\bf u}$ from ${\bf u}^{(0)}$. To do so, we introduce $\Delta_2^-$ on $\ell^2(\mathbb{Z}_-)$ 
with $\mathbb{Z}_- = \{ \dots, -2,-1,0\}$ subject to the Dirichlet condition for $n = 1$. The difference equations (\ref{Eq:u_gen_pq}) for $n \in \mathbb{Z}_-$ can be rewritten in the form 
\begin{equation}
(1 - h(u_n) - \varepsilon \Delta_2^-) u_n = \varepsilon u_1 \delta_{n,0}, \quad n \in \mathbb{Z}_-,
\label{Eq:w}
\end{equation}
where $\delta_{n,0} = 1$ for $n = 0$ and $\delta_{n,0} = 0$ for $n \leq -1$, whereas $h(u) = |u|^{p-1} - \gamma |u|^{q-1}$. Since $p,q \in \mathbb{N}$ and $2 \leq p < q$, there is $C > 0$ such that 
$$
|h(u_n)| \leq C |u_n|, \quad \forall n \in \mathbb{Z}_-,
$$
as long as $\|{\bf u} \|_{\ell^2(\mathbb{Z}_-)} \leq 1$. Since $\|{\bf u} \|_{\ell^2(\mathbb{Z}_-)} = \mathcal{O}(\varepsilon)$ and $u_1 = \mathcal{O}(1)$ as $\varepsilon \to 0$ by the bound (\ref{bound-u}), it follows from (\ref{Eq:w}) that 
$$
u_n = \varepsilon^{|n|+1} u_1^{(0)} \left[ 1 + \mathcal{O}(\varepsilon) \right], \quad n \in \mathbb{Z}_-,
$$
so that the sign of $u_n$ for every $n \in \mathbb{Z}_-$ is the same as the sign of $u_1^{(0)}$ since $\varepsilon > 0$. Similarly, we get 
$$
u_n = \varepsilon^{n-N} u_N^{(0)} \left[ 1 + \mathcal{O}(\varepsilon) \right], \quad n \in \mathbb{Z}_+ = \{ N+1,N+2,\dots\}.
$$
Hence, the sign of $u_n$ for every $n \in \mathbb{Z}_+$ is the same as the sign of $u_N^{(0)}$ since $\varepsilon > 0$. This yields the equality between the number of sign alternations in ${\bf u}$ and that in $\tilde{\bf u}$.
\end{proof}

\end{document}
